% Problem record op_8682c870bea8ab4e. This TeX file is derived from
% database/problems_json/op_8682c870bea8ab4e.json by scripts/sync-tex.mjs; edit the JSON record
% and rerun the script rather than editing this file.
\section{Wigner entropy conjecture}
\label{sec:op_8682c870bea8ab4e}

\paragraph{Problem.}
Does every finite-energy single-mode state with a nonnegative Wigner function have Wigner entropy at least $1+\ln\pi$?
Let $\rho$ be a density operator with $\operatorname{Tr}(\rho a^\dagger a)<\infty$, where $a=(q+ip)/\sqrt2$ and $[q,p]=i$. Require $W_\rho(q,p)\geq0$. Use the normalization $\int_{\mathbb R^2}W_\rho(q,p)\,dq\,dp=1$ and vacuum convention $W_{|0\rangle}(q,p)=\pi^{-1}e^{-q^2-p^2}$. With natural logarithms and $0\ln0=0$, the proposed bound is
\begin{equation}
 h_W(\rho):=-\int_{\mathbb R^2}W_\rho(q,p)\ln W_\rho(q,p)\,dq\,dp
 \geq1+\ln\pi.
 \label{eq:wigner-entropy}
\end{equation}
Prove Eq.~\eqref{eq:wigner-entropy} for all such physical states, or exhibit a density operator violating it.

\subsection*{Status}

Solved

\subsection*{Source}

Van Herstraeten and Cerf state the conjecture in Eq.~(10) of ``Quantum Wigner Entropy'' \sourcecite{ref:wigner-original}{VHC21}. The finite-energy restriction makes the moments and entropies used here finite.

\subsection*{Progress}

\begin{itemize}
  \item Pure Gaussian states attain $h_W=1+\ln\pi$. The bound also holds for passive harmonic-oscillator states, whose Fock-state probabilities decrease with photon number; see Sec.~IV, especially Eq.~(44) \sourcecite{ref:wigner-original}{VHC21}.

  \item Van Herstraeten, Jabbour, and Cerf prove the stronger continuous-majorization relation for Wigner-nonnegative mixtures of the first three Fock states, which implies $h_W\geq1+\ln\pi$ on that family \sourcecite{ref:wigner-majorization}{VJC23}.

  \item Qian and Gagatsos prove the entropy bound for all Wigner-nonnegative states supported on $\operatorname{span}\{|0\rangle,|1\rangle\}$, including coherences (Section II). Section III gives a sufficient condition for a further class of mixed states; it does not cover every Wigner-nonnegative state \sourcecite{ref:wigner-qubits}{QG24}.

  \item Theorem~1 of Van Herstraeten and Cerf proves the bound for outputs obtained by applying a balanced beam splitter to a separable two-mode state and discarding one output. Such outputs are Wigner nonnegative. The theorem is a Wigner--R\'enyi entropy bound for orders at least $1/2$, including the Shannon limit in Eq.~\eqref{eq:wigner-entropy} \sourcecite{ref:wigner-bs}{VHC25}.

  \item Qian and Gagatsos prove the sufficient condition $\mu\leq2/e$, where $\mu=\operatorname{Tr}\rho^2=2\pi\int W_\rho^2$, in Sec.~III.4, Eq.~(26). Their Proposition~1 concerns relaxed phase-space constraints; its extremizers need not represent positive density operators and do not disprove the conjecture \sourcecite{ref:wigner-purity}{QG26}.

  \item He disproves the conjecture with a rank-two state supported on $\operatorname{span}\{|0\rangle,|3\rangle\}$ (Theorem~1, Sec.~5.1). The choice $t=1/1000$ and $\lambda=9/10$ in Eq.~(35) gives $\rho=(1000000|0\rangle\langle0|+|3\rangle\langle3|+900|0\rangle\langle3|+900|3\rangle\langle0|)/1000001$, whose nonzero block has positive determinant $190000/1000001^2$ and whose mean photon number is $3/1000001$. In the convention of Eq.~\eqref{eq:wigner-entropy}, Eqs.~(54)--(55) establish $W_\rho\geq(98199/1000001)W_{|0\rangle}>0$ everywhere and $h_W(\rho)<1+\ln\pi-4.79\times10^{-7}$. The global positivity estimate follows from square completion in Eqs.~(45)--(46), and the strict entropy deficit follows from the analytic polynomial-moment bound in Eqs.~(47)--(53) \sourcecite{ref:wigner-counterexample}{He26}.
\end{itemize}

\subsection*{References}

\begin{enumerate}[label={},leftmargin=4.8em,labelsep=0.5em]
  \footnotesize
  \item[\textup{[VHC21]}]\label{ref:wigner-original}
  Z. Van Herstraeten and N. J. Cerf, ``Quantum Wigner Entropy,'' \emph{Physical Review A} \textbf{104}, 042211 (2021). \href{https://doi.org/10.1103/PhysRevA.104.042211}{doi:10.1103/PhysRevA.104.042211}; \href{https://arxiv.org/abs/2105.12843}{arXiv:2105.12843}.

  \item[\textup{[VHC25]}]\label{ref:wigner-bs}
  Z. Van Herstraeten and N. J. Cerf, ``Wigner Entropy Conjecture and the Interference Formula in Quantum Phase Space,'' \emph{Physical Review A} \textbf{112}, 062207 (2025). \href{https://doi.org/10.1103/1ftk-dm7l}{doi:10.1103/1ftk-dm7l}; \href{https://arxiv.org/abs/2411.05562}{arXiv:2411.05562}.

  \item[\textup{[QG26]}]\label{ref:wigner-purity}
  Q. Qian and C. N. Gagatsos, ``Upper Bounds on the Purity of Wigner Non-negative Quantum States That Verify the Wigner Entropy Conjecture,'' \emph{Physical Review A} \textbf{114}, 022408 (2026). \href{https://doi.org/10.1103/6vcd-qyj9}{doi:10.1103/6vcd-qyj9}; \href{https://arxiv.org/abs/2601.16898}{arXiv:2601.16898}.

  \item[\textup{[VJC23]}]\label{ref:wigner-majorization}
  Z. Van Herstraeten, M. G. Jabbour, and N. J. Cerf, ``Continuous Majorization in Quantum Phase Space,'' \emph{Quantum} \textbf{7}, 1021 (2023). \href{https://doi.org/10.22331/q-2023-05-24-1021}{doi:10.22331/q-2023-05-24-1021}; \href{https://arxiv.org/abs/2108.09167v2}{arXiv:2108.09167v2}.

  \item[\textup{[QG24]}]\label{ref:wigner-qubits}
  Q. Qian and C. N. Gagatsos, ``Wigner Non-negative States That Verify the Wigner Entropy Conjecture,'' \emph{Physical Review A} \textbf{110}, 012228 (2024). \href{https://doi.org/10.1103/PhysRevA.110.012228}{doi:10.1103/PhysRevA.110.012228}; \href{https://arxiv.org/abs/2406.09248}{arXiv:2406.09248}.

  \item[\textup{[He26]}]\label{ref:wigner-counterexample}
  Z. He, ``Physical Counterexamples to the Wigner Shannon Entropy Conjecture,'' preprint (2026). \href{https://arxiv.org/abs/2609.13312v1}{arXiv:2609.13312v1}.
\end{enumerate}

\subsection*{Comment}

The conjecture is false: Theorem~1 of the September 2026 preprint \sourcecite{ref:wigner-counterexample}{He26} supplies a finite-energy density operator with an everywhere positive Wigner function and entropy strictly below the vacuum value, in exactly the normalization used here. Thus the counterexample alternative in the archived statement is settled. The resolving source cited here is a preprint. The previously established bounds for Gaussian, passive, and other restricted classes remain valid; the counterexample is a coherent mixture of the vacuum and three-photon levels and lies outside those proved subclasses.

\subsection*{Field}

Quantum Communication

\subsection*{Topic}

Matrix and entropy inequalities; Gaussian quantum information

\subsection*{ID}

\texttt{op\_8682c870bea8ab4e}
